
\section{Contributions and Acknowledgements}
\label{sec:acknowledgements}

\paragraph{Contributions.}
The project was divided into five main parts.
Raghav Singhal worked on the mathematical foundations and
Model A, Ashwath Chandramouli on Model B and the comparison
of harvesting methods, and Devam Jatin Bhavsar on collapse
thresholds, extinction times and numerical verification.
Mehta Prem worked on the forward PINN analysis, while
Rohan Agrawal worked on inverse PINNs, parameter estimation
and validation. The notebook, report and presentation
were prepared and integrated as a group.

\paragraph{AI assistance.}
We used ChatGPT, OpenAI Codex, Claude Code and
Antigravity AI Assistant during the project.
ChatGPT assisted with mathematical explanations,
review and editing of written material.
OpenAI Codex was used for code development, debugging,
numerical checks, testing and project integration.
Claude Code and Antigravity AI Assistant supported
implementation, numerical experiments and preparation
of code and written drafts, including work related
to the PINNs.

The AI-generated suggestions and outputs were treated
as supporting material. The mathematical results,
numerical findings and submitted explanations remain
the responsibility of the authors.

\paragraph{Sources and Data.}
All population trajectories, reference solutions and
noisy observations used in this project were generated
synthetically from the mathematical models described
in this report. No external population dataset was
used. External mathematical and computational sources
are cited where relevant.

The complete source code, experiment configurations,
numerical data, saved results and figures are available
in the \href{https://github.com/devam2006/mm-allee-harvesting}
{GitHub project repository}.
The accompanying Jupyter notebook reproduces the
main analytical and numerical calculations, executes
short PINN demonstrations and loads the recorded
results of the larger PINN experiments.
